\begin{center}3.\end{center}
\noindent\hspace*{\parindent}Ist eine Reihe \(a',\,a'',\,a'''\) u. s. f. \({.}\,{.}\,{.}\ [R]\), in welcher für jeden endlichen Index das entsprechende Glied einen endlichen reellen Werth erhält, so beschaffen, dass in ihr, so weit man sie auch fortsetzt, kein Glied vorkommt, das grösser\textsuperscript{*}) als eine gegebene Grösse \(\lambda\) wäre, so kann man \(\lambda\) eine obere Grenze für die Reihe nennen (limes supra seriem, une limite en plus); ist hingegen keine Grösse gross genug, um nach diesem Begriffe eine obere Grenze genannt werden zu können, oder mit andern Worten, kann man in der Reihe zu so grossen Grössen als man will oder zu grössern gelangen, wenn man sie nur weit genug fortsetzt, so sagt man die Reihe habe keine obere Grenze. Es ist klar, dass wenn \(\lambda\) eine obere Grenze der Reihe \([R]\) ist, jede Grösse, welche grösser als \(\lambda\) ist, gleichfalls eine obere Grenze der Reihe sein werde, und im Fall nicht schon \(\lambda\) selbst die kleinste obere Grenze ist, wird es noch kleinere als \(\lambda\) geben; nun ist es aber von selbst klar, dass es kleinere Grössen als \(\lambda\) gebe, die nicht mehr obere Grenzen der Reihe genannt werden können; lässt man demnach \(\lambda\) durch alle Zwischengrössen stetig abnehmen, so muss man nothwendig auf eine kleinste obere Grenze \(L'\) kommen, die also die
\par\medskip\noindent\rule{.24\linewidth}{.4pt}\par\smallskip
{\small\noindent\hspace*{\parindent}\textsuperscript{*}) Die Wörter grösser und kleiner werden hier allemal, wo nicht das Gegentheil erinnert wird, mit Rücksicht auf die Zeichen verstanden, so dass 0 und jede positive Grösse grösser als jede negative, und von zwei negativen Grössen diejenige die grössere genannt wird, die ohne Rücksicht auf das Zeichen die kleinere wäre. Eben so sind die Ausdrücke zunehmen und abnehmen zu verstehen.\par}
\clearpage\SourceLeaf{399}\ArithmeticHead{392}{L}{ANALYSIS.\quad NACHLASS.}
\noindent Eigenschaft haben wird, dass kein Glied der Reihe \([R]\) grösser als \(L'\) ist, wohl aber in dieser Reihe, wenn sie nur weit genug fortgesetzt wird, Glieder vorkommen können, die grösser sind, als jede andere Grösse, die kleiner als \(L'\) ist.

\indent Ganz auf ähnliche Weise kann \(\mu\) eine untere Grenze (une limite en moins) der Reihe \([R]\) heissen, wenn in derselben kein Glied vorkommen kann, das kleiner als \(\mu\) wäre, woraus von selbst erhellet, was Reihen sind, die keine untere Grenze haben. Bei jeder Reihe, die untere Grenzen hat, wird es eine grösste untere Grenze \(M'\) geben, so dass jede grössere Grösse als \(M'\) nicht mehr untere Grenze der Reihe heissen kann. — Da wir diese kleinsten obern und grössten untern Grenzen allein brauchen, so wollen wir dieselben in der folge schlechthin die obere, und die untere Grenze nennen und also den vorigen allgemeinern Begriff einer obern oder untern Grenze, welchen wir bloss zur Ableitung des gegenwärtigen gebraucht haben, ganz bei Seite setzen.

\indent Übrigens sieht man leicht ein, dass es zwei verschiedene Arten gebe, wie eine Grösse \(L'\) die obere Grenze einer Reihe sein könne; entweder nemlich, wenn es in der Reihe wirklich ein Glied (oder mehrere) gibt, dass [es] der \(L'\) gleich ist, die übrigen aber alle kleiner sind, oder, wenn zwar kein der Grösse \(L'\) gleiches Glied in der Reihe vorkommen kann, aber doch, wenn \([R]\) weit genug fortgesetzt wird, Glieder, die so wenig als man will von \(L'\) abweichen. Ganz auf dieselbe Art verhält es sich mit den untern Grenzen.

\indent Auf diese Weise gibt es vier verschiedene Arten von Reihen:
\begin{list}{}{\setlength{\leftmargin}{3em}\setlength{\labelwidth}{2.5em}\setlength{\labelsep}{.5em}\setlength{\itemsep}{.35em}\setlength{\parsep}{0pt}}
\item[I.] Reihen, die weder eine obere noch eine untere Grenze haben, z. B. \(1,\allowbreak\,-2,\allowbreak\,+3,\allowbreak\,-4,\allowbreak\,+5\) u. s. f.
\item[II.] Reihen, die keine obere aber eine untere Grenzen haben, wie \(1,\,2,\,3,\,4\) u. s. f., wovon die untere Grenze 1.
\item[III.] Reihen, die keine untere aber eine obere Grenze haben, wie \(-1,\,-2,\,-3\) u. s. f.
\item[IV.] Reihen, die so wohl eine obere als eine untere Grenze haben. Diese beiden Grenzen werden immer ungleich sein, wenn nicht alle Glieder der Reihe einander gleich sind. Z B \(\frac12,\,\frac23,\,\frac34,\,\frac45\) u. s. f , wo die obere Grenze 1, die untere \(\frac12\).
\end{list}
\indent Die letztere Art von Reihen ist hier für uns die wichtigste.

